\title{LongRoPE: Extending LLM Context Window Beyond 2 Million Tokens}

\begin{abstract}
An extensive context window is a sought-after characteristic in large language models (LLMs). Nonetheless, prevailing methods for extending context windows are constrained to approximately 128k tokens, owing to factors such as the significant expense of fine-tuning, the limited availability of sufficiently long textual data, and the destabilizing impact of assigning new values to previously unseen token positions.
\end{abstract}

In this work, we present LongRoPE, a novel method that significantly enlarges the context window of pre-trained LLMs to an outstanding \textbf{2048k} tokens, requiring as few as 1,000 fine-tuning steps with training sequence lengths capped at 256k, all while preserving the performance at the original, shorter context lengths. This advancement is driven by three main contributions: (i) we systematically discover and leverage two distinct non-uniformities within positional interpolation using an efficient search mechanism, facilitating superior initialization for fine-tuning and supporting an eightfold increase in context window without additional fine-tuning; (ii) we propose a staged extension approach, beginning with fine-tuning on a 256k token length and subsequently applying a second round of positional interpolation on this already extended model to reach a 2048k context window; (iii) we further refine LongRoPE at the 8k sequence length to fully restore performance for shorter contexts. Comprehensive empirical studies conducted on LLaMA2 and Mistral across multiple tasks showcase the robustness and practicality of our technique. Notably, LongRoPE-extended models require only minimal adjustments to the positional embedding and otherwise maintain the base model architecture, allowing seamless integration with most existing optimizations. The source code will be made accessible at \texttt{https://github.com/microsoft/LongRoPE}
\end{abstract}

\section{Introduction}

Although Large Language Models (LLMs) have achieved impressive results across a range of applications~\cite{gpt4,llama2}, they are typically constrained by a finite context window, such as the 4096-token capacity of LLaMA2~\cite{llama2}. When processing input that extends beyond this window, model performance deteriorates, as the LLM has not encountered such extended positional sequences during training. This limitation creates significant obstacles in use cases that require processing long input sequences, for example, in-context learning involving multiple examples~\cite{Huang2023FewerIM} and for LLM-based agents~\cite{park2023generative,madaan2023selfrefine}.

Recent studies have demonstrated that it is possible to enlarge the context window of a pre-trained LLM to approximately 128k by applying fine-tuning techniques using longer documents~\cite{longlora,pi,yarn,zhang2024soaring,liu2023scaling}. Nonetheless, three primary challenges impede the extension of context windows beyond this scale. Firstly, the introduction of previously untrained position indices can result in a proliferation of anomalous values, exacerbating out-of-distribution problems and hindering the convergence of the fine-tuning process~\cite{pi}. This difficulty becomes especially pronounced when expanding from 4k to over $>$1000k, a scenario in which more than 90\% of positions are entirely new. Secondly, successful fine-tuning generally depends on access to sufficiently long documents that match the targeted context window lengths. Current datasets contain a scarcity of such ultra-long texts, particularly for lengths beyond 1000k. In addition, processing these lengthy inputs is extremely resource-intensive, necessitating considerable GPU power and extended training durations. Thirdly, increasing the context window to extreme lengths causes the model’s attention to be distributed over an immense number of token positions, resulting in a dilution of focus that impairs performance on tasks involving shorter contexts~\cite{pi}.

A common strategy to address the initial challenge involves interpolating RoPE positional embeddings~\cite{rope,pi}, wherein newly encountered position indices are rescaled to align with the original pretraining interval, as illustrated in Fig.~\ref{fig:overview}. Position Interpolation (PI)~\cite{pi} achieves this by adjusting RoPE’s rotary angles through a linear interpolation guided by the extension ratio. NTK~\cite{ntk,dynamicntk} introduces an alternative approach, recommending a mixture of non-uniform interpolation and extrapolation methods over different RoPE dimensions. YaRN~\cite{yarn} further expands on this idea by dividing RoPE’s dimensions into three groups according to frequency, and separately applying extrapolation, NTK, and linear interpolation strategies to each group. Nonetheless, positional embeddings within Transformer models possess a \emph{complex, non-uniform distribution of information entropy}. Existing techniques tend to overlook this nuanced non-uniformity, thereby causing a loss of information and restricting the ability to expand the context window size effectively.

Section~\ref{sec:analysis} presents empirical evidence for two primary conclusions: \textbf{(1)} Successful positional interpolation must account for two distinct sources of non-uniformity: differences in RoPE dimensions as well as variations in token positions. It is advantageous to apply less interpolation to lower RoPE dimensions and to tokens near the start of the sequence; however, the optimal interpolation strategy is contingent on the specific extended length being targeted. 
  \textbf{(2)} Incorporating these sources of non-uniformity into the positional interpolation procedure allows the model to better preserve key information encoded by the original RoPE, particularly for crucial dimensions and token indices. This strategy reduces information loss resulting from interpolation, thereby leading to superior initialization during fine-tuning. In addition, it enables up to an 8$\times$ increase in sequence length without requiring fine-tuning.

Building on these observations, we present \textbf{LongRoPE}, a robust technique designed to push the LLM context window past 2 \emph{million} tokens. LongRoPE incorporates three principal innovations. To begin with, {LongRoPE} leverages the multidimensional, non-uniform characteristics of positional interpolation to the fullest. It determines optimal rescale factors for RoPE's rotation angles in each dimensional component, contingent on token positions. Given that the space for identifying these rescale factors grows exponentially with the desired extension ratio, {LongRoPE} employs an evolutionary search algorithm enhanced by two optimization strategies that markedly improve the efficiency of the search. An illustration of the rescaled RoPE produced by this search process can be found in Fig.~\ref{fig:overview}.

Subsequently, {LongRoPE} employs an efficient, step-wise extension approach to realize a 2048k context window, circumventing the necessity for direct fine-tuning on extremely lengthy texts, which are both rare and difficult to obtain. Initially, the method identifies a 256k context length supported by the pre-trained LLM and applies fine-tuning at this scale. Owing to the capacity of our non-uniform positional interpolation to facilitate an 8$\times$ expansion even in the absence of further fine-tuning, a subsequent search is carried out to determine new RoPE rescale factors using the fine-tuned, length-extended LLM. This process enables {LongRoPE} to attain a 2048k context window for LLaMA2 and Mistral~\cite{mistral}.

Lastly, to address potential drops in performance when handling the initial (smaller) context window, {LongRoPE} further refines the RoPE rescale factors in the expanded LLM. In parallel with the process of increasing the context length from 256k to 2048k, we also apply our search algorithm to reduce context windows to 4k and 8k on the LLM that has been fine-tuned at 256k. This aims to minimize positional interpolation effects. For inference scenarios where the input sequence length falls below 8k, we revise the RoPE settings using the rescale factors determined through our search procedure.

Through comprehensive experimentation involving multiple LLMs and a range of long-context tasks, we validate the efficacy of our proposed approach. Our results indicate that LongRoPE consistently preserves low perplexity across evaluation lengths spanning from 4k up to 2048k, achieves passkey retrieval accuracies exceeding 90\%, and attains accuracy on par with baselines in established benchmarks set within a 4096-token context window. Additionally, LongRoPE is compatible with all LLMs utilizing RoPE embedding. The source code and LongRoPE-2048k model checkpoints will be made publicly available.

\section{Non-uniformity in Positional Interpolation}
\label{sec:analysis}

\subsection{Preliminary}

Transformer architectures depend on the incorporation of explicit positional cues, typically achieved via position embeddings, to capture the sequential order of input tokens. In this study, we concentrate on RoPE~\cite{rope} position embedding, a method that has become prevalent in state-of-the-art LLMs. Given a token located at position $n$, its associated RoPE encoding can be succinctly represented as:

\begin{equation}
		\fontsize{7.8}{7.8} \selectfont
	\label{eq:rope}
	\left[ cos(n\theta_0), sin(n\theta_0), cos(n\theta_1), \cdot\cdot\cdot, cos(n\theta_{d/2-1}), sin(n\theta_{d/2-1}) \right]   
\end{equation}

where $d$ denotes the dimension of the embedding space, 
$n\theta_i$ corresponds to the rotary angle assigned to the token situated at position $n$, 
and $\theta_i=\theta^{-2i/d}$ specifies the rotation frequencies. Within RoPE, the typical base value utilized for $\theta$ is 10000.

\textbf{\emph{Context window extension ratio $s$} and positional interpolation}. We define $s$ as the quotient of the extended context length $L'$ over the initial context length $L$: $s=\frac{L'}{L}$.

To increase the context window from $L$ to $L'$, existing positional interpolation techniques propose modifying the rotation frequencies $\theta_i$ by scaling them according to the extension ratio $s$. If we define $\beta=\theta^{2/d}$ and use $\lambda$ to represent the true scaling factor corresponding to $s$, then these various positional interpolation methods can be generalized in the following way:

\begin{equation}	
	\fontsize{7.5}{7.5} \selectfont
	\label{eq:unified_rope}
	\begin{aligned}
		\left[cos\left(\frac{n}{\lambda(\beta)^0}\right), sin \left(\frac{n}{\lambda(\beta)^0}\right), cos\left(\frac{n}{\lambda(\beta)^1}\right), \cdot\cdot\cdot,  sin \left(\frac{n}{\lambda(\beta)^{d/2-1}}\right) \right]   
	\end{aligned}
\end{equation}

\textbf{Linear positional interpolation (PI)}. PI~\cite{pi} employs a method where position indices are linearly interpolated, constrained within the maximum sequence length encountered during pre-training. When the sequence is extended by a factor $s$, the rotation angles corresponding to each position are uniformly scaled down by $\lambda=s$ across every RoPE dimension. Despite this approach, such compression of positional representation leads to densely packed positional cues, which impairs the model's capacity to differentiate tokens that are near each other in the sequence. Consequently, PI's effectiveness diminishes as the extension ratio increases.

\textbf{NTK-based interpolation and extrapolation}. The works of ~\cite{ntk,dynamicntk} analyze RoPE from the standpoint of information encoding and incorporate insights from Neural Tangent Kernel (NTK) theory~\cite{jacot2018neural,tancik2020fourier}. In response to the dense clustering of positions in PI, these methods propose redistributing the interpolation load among the various RoPE dimensions. Specifically, they reduce the scaling applied to dimensions corresponding to higher frequencies, while amplifying it for those representing lower frequencies. This approach yields improvements in both positional interpolation and extrapolation, formalized with $\lambda=s^{i}$. The dynamic NTK technique~\cite{dynamicntk} further refines this by adapting the extension ratio dynamically for each position, depending on the current sequence length. Contrary to PI, which requires additional fine-tuning, NTK-based strategies allow for the expansion of context windows even without fine-tuning, though this typically caps the maximum achievable extension at 4$\times$.

\textbf{YaRN}~\cite{yarn} presents a notable enhancement in the effectiveness of positional interpolation. This method segments RoPE dimensions into three distinct groups according to their frequency characteristics, applying a tailored interpolation approach for each. Specifically, extrapolation ($\lambda$=1) is applied to dimensions associated with higher frequencies, while lower frequency dimensions make use of linear interpolation (PI). For those RoPE dimensions whose frequencies fall in the intermediate range, the NTK technique is utilized. The central innovation of YaRN is its categorization of RoPE dimensions; however, this relies primarily on empirical grouping by researchers, which could lead to less than optimal outcomes when adapting to newly developed LLMs.

\subsection{Study on Non-uniform Positional Interpolation}

Building on the insights from NTK and YaRN, we observe that their improvements stem from introducing non-linear mechanisms, especially through handling varying frequencies in RoPE dimensions for tailored interpolation and extrapolation. Yet, prevailing non-linear functions are largely based on hand-crafted approaches. This observation prompts two key questions: (1) Does existing positional interpolation represent the best possible solution? (2) Might there be additional non-linear strategies yet to be discovered?

In order to address these issues, we employ evolution search (refer to Sec.~\ref{sec:method}) to identify improved non-uniform positional interpolations for LLaMA2-7B. The search process is steered by perplexity metrics, evaluated on 5 randomly selected samples from the PG19~\cite{pg19} validation dataset. Our experimental study yields several notable observations.

\textbf{Finding 1}: \textit{RoPE dimensions exhibit substantial non-uniformities, which are not effectively handled by current positional interpolation methods.}

We optimize $\lambda$ independently for each RoPE dimension as specified in Eq.~\ref{eq:unified_rope}. In Table~\ref{tbl:exp1}, we present a comparison of the perplexity results obtained with various methods applied to LLaMA2-7B, evaluated on the PG19 and Proof-pile~\cite{proof-pile} test sets, all without fine-tuning. Our approach achieves notable perplexity reductions, indicating that prevalent linear (PI) and non-uniform (Dynamic-NTK and YaRN) interpolation strategies are not optimal. Importantly, YaRN performs worse than both PI and NTK on the PG19 benchmark, as it fails to accommodate the intended context window length for large language models that have not been fine-tuned. Specifically, the perplexity of YaRN sharply increases beyond 7k tokens in the case of an 8k context window.

In our investigation, we observe that the rescaling factors $\lambda$ in Eq.~\ref{eq:unified_rope} are not uniform; this stands in contrast to the consistent scale $s$ applied in the formulations of PI, NTK, and YaRN's group-based approach. These varying rescaling factors yield marked enhancements in the language modeling abilities of LLaMA2, as measured by perplexity, within 8k and 16k context lengths—accomplished without any model fine-tuning. The reason for this improvement is that the resulting positional encodings more faithfully maintain the characteristics of the original RoPE, particularly across key dimensions, thereby alleviating the challenge for the LLM in differentiating between tokens that are positioned closely together.

\textbf{Finding 2}: \textit{RoPE for the initial tokens in the input sequence should be extrapolated with less interpolation.} 

We propose that the RoPE of the first $\hat{n}$ tokens within input sequences should involve minimal interpolation. This hypothesis arises from their prominent attention scores, rendering them vital in the computation of attention within the layers, a phenomenon noted in Streaming LLM~\cite{streamingllm} and LM-Infinite~\cite{han2023lminfinite}. To examine this hypothesis, we enlarge the context window to 8k and 16k using PI and NTK, selectively exempting the first $\hat n$ tokens (with values of 0, 2, ..., 256) from interpolation. Setting $\hat n$=0 corresponds to the standard PI and NTK methods. Table~\ref{tbl:tokenposition} demonstrates two principal findings: \textbf{(1)} excluding the initial tokens from position interpolation consistently enhances model performance; \textbf{(2)} the most effective choice of $\hat n$ varies according to the target length for context extension.

\textbf{Finding 3}: \textit{Non-uniform positional interpolation effectively extends LLM context window in both fine-tuning and non-fine-tuning settings.} 

Although our experiments demonstrate that the proposed non-uniform position interpolation method yields substantial performance gains for extending context lengths to 8k and 16k tokens without fine-tuning, achieving robust results at even larger context lengths necessitates fine-tuning. Therefore, we perform fine-tuning on LLaMA2-7B using our identified RoPE modifications for a 64k token context window (refer to Appendix for detailed setup). As illustrated in Table~\ref{tbl:finetune}, our approach surpasses both PI and YaRN by a notable margin, regardless of whether fine-tuning has been applied. This superiority stems from more effectively leveraging non-uniform positional interpolation, which preserves information more efficiently and offers superior parameter initialization for fine-tuning.

\textbf{Summary}. This work identifies two non-uniform patterns—variations in RoPE dimensions and across token positions—and demonstrates that harnessing these non-uniformities in positional interpolation substantially enhances context window extension for LLMs.

\section{LongRoPE}

\label{sec:method}
Building on these insights, we propose LongRoPE, a method that initially applies an efficient search algorithm designed to leverage both forms of non-uniformity. This approach is subsequently employed to expand the context window of LLMs to more than 2 million tokens.

\subsection{Problem Formulation}

These two forms of non-uniformity greatly expand the solution space, thereby complicating the optimization process. To manage this, we recast the multidimensional non-uniform position interpolation optimization challenge as a search problem.

Considering an LLM that supports a context window of size $L'$ and processes input documents $\mathbf{X}$ composed of sequences $\mathbf{x}\in\mathbf{X}$ where each sequence's token count exceeds $L'$, we express the original rotary angle for the $i^{th}$ dimension in the RoPE embedding at token position $n$ as $\frac{n}{\beta^i}$. The corresponding optimization task is represented as:

\begin{equation}
\begin{gathered}
		\label{eq:problem}

\underset {\mathbf{x}\in\mathbf{X}; \, |\mathbf{x}|\ge L'}{ \operatorname {arg\,min} } \,  \mathcal{L}\, \left( \text{LLM} (\text{RoPE}, \mathbf{X})\right), \text{where} 
\\
\scalebox{0.93}{
$\underset {\begin{subarray}{l} i=0,\cdot\cdot,\frac{d}{2}-1;\\ n\in[ 0,|\mathbf{x}|);\end{subarray}}{\text{RoPE}_(n)}= \left[\cdots, \cos\left(\mathbb{I}(\hat{\lambda}_{i}, \hat{n})\frac{n}{\beta^i}\right), \sin\left(\mathbb{I}(\hat{\lambda}_{i}, \hat{n})\frac{n}{\beta^i}\right), \cdots\right] $}\\
\text{with} \; \mathbb{I}(\hat{\lambda}_{i},\hat{n}) = \begin{cases}
1 & \text{if $n < \hat{n}$}\\
\frac{1}{\lambda_{i}} & \text{if $n \geq \hat{n}$}
\end{cases}
\end{gathered}
\end{equation}

In this context, we define a set of rescaling factors, $\mathbb{I}(\hat{\lambda}_{i},\hat{n})$, designed to address both types of non-uniformity present in the model. Here, $\hat{\lambda_i}$ captures the non-uniformity specific to RoPE dimensions, while $\hat{n}$ represents the threshold for token positions. The rescaling function $\mathbb{I}(\hat{\lambda}_{i},\hat{n})$ is utilized to adjust the rotational angle in the $i^{th}$ RoPE dimension; $\hat\lambda_{i}$ functions as the scaling coefficient, and $\hat{n}$ as the positional cutoff. For token positions prior to $\hat{n}$, that is, for the first $\hat{n}$-1 tokens, the rescaling factor $\hat\lambda_{i}$ is not applied, and the standard RoPE rotation angle $\frac{n}{\beta^i}$ remains in use. For token positions where $n \ge \hat{n}$, the rescale factor is incorporated into the rotation calculation.

Assuming a desired context window length of $L'$, the task involves identifying the optimal set of rescale factors ($\mathbb{I}(\hat{\lambda}_{0},\hat{n})$, $\mathbb{I}(\hat{\lambda}_{1},\hat{n})$, ..., $\mathbb{I}(\hat{\lambda}_{i},\hat{n})$, ...) corresponding to each RoPE dimension from the first through the $d^{th}$. By applying these adjusted rescaling parameters to the target $\text{LLM}$'s $\text{RoPE}$, the model is expected to minimize the next-token prediction loss, $\mathcal{L}$ (equivalent to the perplexity), for any input sequence $\mathbf{X}$ whose length equals $L'$.

\subsection{Searching the Non-uniform Position Interpolation}

To address the problem stated in Eq.~\ref{eq:problem}, we present a straightforward yet powerful approach that identifies optimal RoPE rescale factors, enabling the model to take full advantage of the multidimensional variability present in positional embeddings.

\textbf{Search space}. We construct an extensive search space designed to accommodate the two identified non-uniformities. The structure of this search space is depicted in Table~\ref{tbl:searchspace}. Specifically, our method permits the determination of a unique rescale factor for every individual dimension within the RoPE embedding. To streamline the design of the search space, the algorithm explores $\lambda_i$ and $\hat{n}$, rather than directly searching for $\mathbb{I}(\hat{\lambda}_{i},\hat{n})$, where we set $\hat{\lambda}_{i}=1/\lambda_i$. According to Table~\ref{tbl:searchspace}, $\lambda_i$ may be selected from a range that starts at 1.0 (representing standard extrapolation) up to $s\times1.25$ (indicating broader interpolation than PI), incremented in steps of 0.01, where $s$ corresponds to the intended expansion ratio for the context window.

$\hat{n}$ determines how many of the earliest token positions preserve the standard RoPE embedding, foregoing position interpolation. In practice, we permit $\hat{n}$ to take values from the set \{0, 1, 2, 4, 8, 12, 16, 20, 24, 28, 32, 64, 128, 256\} during the search. Setting $\hat{n}=0$ results in every token position adopting the learned rescale factors.

\textbf{Evolution-based search}. The search space defined in Table~\ref{tbl:searchspace} encompasses a vast array of positional interpolation configurations, making it difficult to explore efficiently. For instance, when employing a $s=4\times$ extension, the search complexity becomes $400^{128/2}\times14=4\times10^{167}$ possible options. As the extension ratio increases, the size of the search space grows exponentially. To tackle this, we employ evolution search~\cite{guo2020single}, and further incorporate two optimization strategies to significantly enhance the efficiency of the search process. The overall search method is detailed in Algorithm~\ref{alg:search}.

\textit{Optimized initial population generation}. In lieu of randomly generating all $P$ rescale factors to form the initial population, we explicitly include the three RoPE rescale factors associated with PI, NTK, and YaRN as dedicated individuals. The other $P$-3 members of the population are then produced by randomly applying mutations to these three rescale factors with probability $p$.

\textit{Monotonically non-decreasing constraint}. Following the creation of the initial population, we evaluate each individual by calculating LLM perplexity. This involves applying the associated RoPE rescale factors to the target LLM and then measuring the perplexity on the input $\mathbf{X}$. The best-performing $k$ individuals are then selected as parents for the subsequent evolutionary steps. Yet, due to the large search space, straightforward mutation and crossover operations may frequently yield suboptimal candidates, resulting in superfluous perplexity evaluations. This inefficiency becomes especially pronounced for larger values of $L'$, since each perplexity assessment incurs substantial computational cost.

To mitigate this issue, we enforce a monotonic non-decreasing constraint on the RoPE rescaled factors being sampled, such that $\lambda_i \le \lambda_{i+1}$. Only those RoPE configurations meeting this criterion are considered during perplexity assessments for LLMs, resulting in a substantial reduction of the search space. More precisely, we stipulate that $\lambda_i$ must increase monotonically with respect to the RoPE dimension index ($i=0,\ldots,63$). The motivation for imposing this dimension-wise monotonicity derives from NTK theory~\cite{jacot2018neural,tancik2020fourier,ntk}, which indicates that lower-dimensional components associated with higher frequencies necessitate reduced interpolation (implying smaller $\lambda_i$ values), whereas higher-dimensional, lower-frequency components are capable of tolerating greater interpolation (allowing for larger $\lambda_i$ values).

\begin{algorithm}[t]
	\small
	\caption{The search algorithm}
	\textbf{Input:} target LLM, input samples $\mathbf{X}$, population size $P$, mutation size $N_1$, crossover size $N_2$, max iterations $\mathcal{T}$, mutate probability $p$\\

	\begin{algorithmic}[1]
		\label{alg:search}
		\STATE Initialize Top-k as empty;	
		\STATE Set initial population $\text{P}_0$ using \textit{Initialize\_population\_with\_optimization} ($P$, $\mathbf{X}$, $p$); 
		\FOR{$i$ from $1$ to $\mathcal{T}$}
		\STATE \textit{Compute\_perplexity} (LLM, $\text{P}_{i-1}$, $\mathbf{X}$);
		\STATE  Update Top-k via \textit{Update\_Topk} (Top-k, $\text{P}_{i-1}$);
		\STATE Generate mutated candidates: $\text{P}_{mutation}$ = \textit{Mutation\_with\_mono\_constraint} (Top-k, $N_1$, $p$); 
		\STATE Generate crossover candidates: $\text{P}_{crossover}$ = \textit{Crossover\_with\_mono\_constraint} (Top-k, $N_2$);
		\STATE Combine for next generation: $\text{P}_i$ = $\text{P}_{mutation}$ $\cup$ $\text{P}_{crossover}$ $\cup$ Top-k;
		\ENDFOR
		\STATE Output the member with minimum perplexity found in Top-k;
	\end{algorithmic}
\end{algorithm}

\textbf{8$\times$ extension without fine-tuning}. Through evolutionary search, our approach successfully uncovers optimal non-uniform RoPE rescaling parameters, selectively maintaining essential dimensions and positional encodings to reduce information degradation due to interpolation. As illustrated in Fig.\ref{fig:non-ft}, this strategy enables LLaMA2’s context window to expand from 4k up to 32k tokens without any need for fine-tuning. Conversely, alternative techniques like PI, as well as non-uniform NTK and YaRN, experience a sharp increase in perplexity once the context length is doubled.

\subsection{Extending LLM Context Window to 2048K}
\label{sec:extendto2048k}

\textbf{Progressive extension to 2048k}. In this section, we present our approach for expanding the context window of pre-trained LLMs from the standard 4k up to 2048k tokens. As previously shown, applying our non-uniform positional interpolation enables an 8$\times$ increase in context length without any fine-tuning. However, achieving much larger extensions, such as 512$\times$, necessitates fine-tuning the model. One possible strategy involves identifying suitable RoPE rescale factors for the target context size of 2048k, followed by fine-tuning. Nonetheless, this approach is hindered by the extremely high computational demands of such training. Additionally, our empirical observations indicate that effectively fine-tuning LLMs for such large extension ratios poses significant difficulties (refer to Appendix).

Fortunately, LongRoPE demonstrates efficacy for both base and fine-tuned extended LLMs. To this end, we propose an efficient, staged approach that reaches the desired 2048k context length using only 1k fine-tuning steps, all conducted within a training length capped at 256k.

$\Diamond$ \textit{LongRoPE search for expanding pre-trained LLMs to 256k}. Using LLaMA2 for illustration, we perform a search aimed at context window sizes of 128k and 256k. Here, the extension ratios correspond to 32$\times$ and 64$\times$, respectively.

$\Diamond$ \textit{Fine-tuning to 256k}. After pre-training, the LLM undergoes additional fine-tuning to extend the context window to 256k. The procedure starts by fine-tuning LLaMA2 for 400 steps utilizing RoPE rescaled factors set for 128k. Subsequently, these rescaled factors are swapped out for their 256k counterparts at the obtained checkpoint, followed by a further 600 steps of fine-tuning. This staged strategy has shown to be more effective compared to directly fine-tuning the model to 256k.

$\Diamond$ \textit{Expanding fine-tuned extended LLM to 2048k via LongRoPE search}. In the last step, we carry out an additional search process on the LLM that was previously fine-tuned to handle 256k token sequences. This approach enables us to achieve a remarkable context window length of 2048k, all without any extra fine-tuning. As a result, the context window is expanded by a factor of $512\times$.

\textbf{Recovery of performance in shorter context windows}. Upon increasing the context window size to an extensive 2048k, we observe a decline in effectiveness for sequences situated within the initial context window. This phenomenon is well-documented in the literature on positional interpolation~\cite{pi}, where mapping position embeddings to a higher-dimensional space constrains them to a compact area in the original window, thereby diminishing the language model’s capability. Specifically, when utilizing a 512$\times$ extension, positions in the original 4k context window become densely packed, exacerbating this adverse effect.

To address this issue, an additional evolution search is conducted on the expanded LLM to fine-tune the RoPE rescale factors specifically for shorter context lengths (such as 4k and 8k). The upper bound for the searched $\lambda$ is lowered since shorter sequences necessitate less positional interpolation. At inference time, the LLM adaptively modifies the related RoPE rescale factors.

\section{LongRoPE}

\label{sec:method}
Building upon these insights, we propose LongRoPE. This approach initially devises an effective search algorithm designed to leverage the identified non-uniformities, and subsequently applies this method to expand the LLM context window to support more than 2 million tokens.

\subsection{Problem Formulation}

The presence of these two non-uniformities expands the range of potential solutions and adds layers of difficulty to the optimization process. To manage this challenge, we recast the multidimensional non-uniform position interpolation optimization task as a search problem.

Consider a large language model (LLM) designed to accommodate a context window of size $L'$, while processing extended input sequences $\mathbf{X}$ such that each element $\mathbf{x}\in\mathbf{X}$ contains more tokens than $L'$. For the $i^{th}$ dimension in the RoPE embedding, the canonical rotary angle at token position $n$ is represented by $\frac{n}{\beta^i}$. The optimization task can thus be described as follows:

\begin{equation}
\begin{gathered}
		\label{eq:problem}

\underset {\mathbf{x}\in\mathbf{X}; \, |\mathbf{x}|\ge L'}{ \operatorname {arg\,min} } \,  \mathcal{L}\, \left( \text{LLM} (\text{RoPE}, \mathbf{X})\right), \text{with \;} 
\\
\scalebox{0.93}{
$\underset {\begin{subarray}{l} i=0,\cdot\cdot,\frac{d}{2}-1;\\ n\in[ 0,|\mathbf{x}|);\end{subarray}}{\text{RoPE}_(n)}= {\left[\cdot\cdot,\cos\left(\mathbb{I}(\hat{\lambda}_{i}, \hat{n})\cdot\frac{n}{\beta^i}\right),  \sin\left(\mathbb{I}(\hat{\lambda}_{i}, \hat{n})\cdot\frac{n}{\beta^i}\right),\cdot\cdot\right] }$}\\
	\text{with \;} \mathbb{I}(\hat{\lambda}_{i},\hat{n})=\begin{cases}
	1&\mbox{\;  $n<\hat{n}$}\\
{\frac{1}{\lambda}_{i}}&\mbox{\; $n\geq\hat{n}$}
\end{cases}
	\end{gathered}
\end{equation}

In this approach, we define a group of rescaling factors, denoted as $\mathbb{I}(\hat{\lambda}_{i},\hat{n})$, to address both types of non-uniformities. Here, $\hat{\lambda_i}$ represents the dimension-specific non-uniformity in RoPE, while $\hat{n}$ indicates non-uniformity concerning token positions. The function $\mathbb{I}(\hat{\lambda}_{i},\hat{n})$ is employed to adjust the rotational angle associated with the $i^{th}$ dimension of RoPE, utilizing $\hat{\lambda}_i$ as the scaling parameter and $\hat{n}$ as the threshold for token positions. For token positions less than $\hat{n}$, that is, for the first $\hat{n}$-1 tokens, the rescaling factor $\hat{\lambda}_i$ is not applied and the conventional RoPE rotational angle $\frac{n}{\beta^i}$ remains in use. When the position index satisfies $n \geq \hat{n}$, the rescaling factor is incorporated.

With a specified target context window of size $L'$, our aim is to determine the most effective rescale factors ($\mathbb{I}(\hat{\lambda}_{0},\hat{n})$, $\mathbb{I}(\hat{\lambda}_{1},\hat{n})$, ... $\mathbb{I}(\hat{\lambda}_{i},\hat{n})$, etc.) corresponding to each RoPE dimension from the $1^{st}$ through the $d^{th}$. By applying these rescaling coefficients to the $\text{RoPE}$ of the target $\text{LLM}$, we seek to minimize the next-token prediction loss $\mathcal{L}$ (which corresponds to the model’s perplexity) for test sequences $\mathbf{X}$ whose length is $L'$.

\subsection{Searching the Non-uniform Position Interpolation}

To address the challenge outlined in Eq.~\ref{eq:problem}, we propose a straightforward yet robust approach that identifies optimal RoPE rescale factors, thereby leveraging the full extent of multidimensional non-uniformities present within position embeddings.

\textbf{Search space}. Our method constructs an expansive search space capable of capturing both identified non-uniformities, as summarized in Table~\ref{tbl:searchspace}. Specifically, we enable the selection of unique rescale factors for every individual dimension within the RoPE embedding. To streamline the formulation of the search space, we opt to search for $\lambda_i$ and $\hat{n}$ directly, instead of operating on $\mathbb{I}(\hat{\lambda}_{i},\hat{n})$, where the relationship $\hat{\lambda}_{i}=1/\lambda_i$ holds. As detailed in Table~\ref{tbl:searchspace}, $\lambda_i$ may be chosen from a range that begins at 1.0 (corresponding to direct extrapolation) and extends up to $s\times1.25$ (corresponding to interpolation that exceeds that of PI), incrementing in steps of 0.01. Here, $s$ denotes the intended ratio by which the context window is to be expanded.

The parameter $\hat{n}$ determines how many leading token positions maintain their original RoPE embedding, thereby bypassing position interpolation. In practice, we explore values of $\hat{n}$ drawn from the set \{0, 1, 2, 4, 8, 12, 16, 20, 24, 28, 32, 64, 128, 256\}. Setting $\hat{n}=0$ means that every token position employs the optimized rescale factors.

\textbf{Evolution-based search}. The search space described in Table~\ref{tbl:searchspace} encompasses a vast array of positional interpolation configurations, making efficient navigation highly nontrivial. For instance, with an extension factor of $s=4\times$, the total number of possibilities reaches $400^{128/2}\times14 = 4\times10^{167}$. As the extension ratio increases, the search space grows at an exponential rate. To tackle this complexity, we adopt evolution search~\cite{guo2020single}, incorporating two optimization strategies to significantly enhance the efficiency of the search process. The overall methodology is outlined in Algorithm~\ref{alg:search}.

\textit{Optimized initial population generation}. In place of assigning all $P$ rescale factors at random during initialization, we explicitly include the three RoPE rescale factors representing PI, NTK, and YaRN as members of the initial population. The other $P$-3 individuals are produced by applying random mutations to these three rescale factors, with each mutation occurring with probability $p$.

\textit{Monotonically non-decreasing constraint}. Once the initial population is established, we evaluate the LLM perplexity for every candidate. This involves using each individual’s respective RoPE rescale factors within the target LLM to assess perplexity on the input $\mathbf{X}$. The best-performing $k$ candidates are selected as parents for the evolutionary phase. Nevertheless, due to the immense search space, simple mutation and crossover strategies often traverse unpromising regions, resulting in superfluous perplexity evaluations. This inefficiency is exacerbated as $L'$ grows, since each perplexity computation is resource-intensive and time-consuming.

To tackle this issue, we enforce a monotonic non-decreasing restriction on the RoPE rescaled factors sampled, such that $\lambda_i\le \lambda_{i+1}$. Only those RoPE settings fulfilling this requirement are considered during LLM perplexity assessments, thereby considerably decreasing the computational effort required for searching. Concretely, we stipulate that $\lambda_i$ grows monotonically along with the RoPE dimension (for $i=0,...,63$). This constraint follows the rationale from NTK theory~\cite{jacot2018neural,tancik2020fourier,ntk}, which implies that dimensions associated with higher frequencies (lower indices) benefit from reduced interpolation (hence, smaller $\lambda_i$), whereas those with lower frequencies (higher indices) allow for increased interpolation (thus, larger $\lambda_i$).

\begin{algorithm}[t]
	\small
	\caption{The search algorithm}
	\textbf{Input:} target LLM, input samples $\mathbf{X}$, population size $P$, mutation size $N_1$, crossover size $N_2$, maximum iterations $\mathcal{T}$, mutation probability $p$\\

	\begin{algorithmic}[1]
		\label{alg:search}
		\STATE Top-k $\leftarrow$ $\phi$;	
		\STATE $\text{P}_0 \leftarrow$ \textit{Initialize\_population\_with\_optimization} ($P$, $\mathbf{X}$, $p$); 
		\FOR{$i = 1$ \TO $\mathcal{T}$}
		\STATE \textit{Compute\_perplexity} (LLM, $\text{P}_{i-1}$, $\mathbf{X}$);
		\STATE Top-k $\leftarrow$ \textit{Update\_Topk} (Top-k, $\text{P}_{i-1}$);
		\STATE $\text{P}_{mutation}$ $\leftarrow$ \textit{Mutation\_with\_mono\_constraint} (Top-k, $N_1$, $p$); 
		\STATE $\text{P}_{crossover}$ $\leftarrow$ \textit{Crossover\_with\_mono\_constraint} (Top-k, $N_2$);
		\STATE $\text{P}_i$ $\leftarrow$ $\text{P}_{mutation} \cup \text{P}_{crossover} \cup$ Top-k;
		\ENDFOR
		\STATE Return the individual within Top-k exhibiting the minimal perplexity;
	\end{algorithmic}
\end{algorithm}

\textbf{8$\times$ extension without fine-tuning}. Through evolutionary search, we successfully discover non-uniform RoPE rescale parameters that safeguard essential dimensions and token positions, thereby reducing information degradation due to interpolation. As shown in Fig.\ref{fig:non-ft}, our approach allows LLaMA2’s context window to be expanded from 4k to 32k tokens without any fine-tuning. By comparison, alternative techniques like PI and the non-uniform versions of NTK and YaRN exhibit a marked increase in perplexity beyond a 2$\times$ window extension.

\subsection{Extending LLM Context Window to 2048K}
\label{sec:extendto2048k}

\textbf{Progressive extension to 2048k}. In this section, we present our approach for scaling the context window of pre-trained LLMs from the conventional 4k up to beyond 2048k tokens. Our non-uniform positional interpolation is capable of providing an 8$\times$ increase in context length without requiring any fine-tuning. However, when even greater extensions are desired (e.g., up to 512$\times$), fine-tuning becomes indispensable. One possible strategy involves identifying appropriate RoPE rescaling parameters for the target window size of 2048k, followed by a fine-tuning stage. Nevertheless, this process is hindered by the substantial computational costs associated with such large-scale training. Additionally, our practical experience suggests that effectively fine-tuning LLMs under such significant extension ratios presents considerable difficulties (refer to Appendix).

Fortunately, LongRoPE demonstrates efficacy with both the base model and the extended LLM after fine-tuning. Consequently, we propose a streamlined, incremental approach that reaches the desired 2048k context length through only 1k fine-tuning steps, all conducted at a training sequence length no greater than 256k.

$\Diamond$ \textit{LongRoPE search for adapting pre-trained LLMs to 256k}. Using LLaMA2 for illustration, we perform a search targeting context window sizes of 128k and 256k. This process results in an extension ratio of 32$\times$ and 64$\times$, respectively.

$\Diamond$ \textit{Fine-tuning to 256k}. Next, we adapt the pre-trained LLM to accommodate a 256k context window through additional fine-tuning. Initially, we conduct 400 fine-tuning steps on LLaMA2 utilizing RoPE rescaled factors tailored for 128k. After this stage, we update the RoPE rescaled factors to correspond to 256k within the obtained checkpoint, followed by a further 600 fine-tuning steps. Compared to direct fine-tuning with a 256k context window, this staged approach is notably more efficient.

$\Diamond$ \textit{Expanding fine-tuned extended LLM to 2048k via LongRoPE search}. In the last step, an additional search is conducted on the 256k-length LLM that has already been fine-tuned. Through this process, the context window is successfully enlarged to 2048k, all without the need for additional fine-tuning. The achieved extension ratio is $512\times$.

\textbf{Restoring performance in shorter context windows}. Upon increasing the context window to a notably long 2048k tokens, we observe that performance within the standard context window declines. This phenomenon is well-documented with positional interpolation methods~\cite{pi}, which constrain high-dimensional position embeddings inside the original context window to a significantly compressed space, thereby impairing the model’s effectiveness. Specifically, with an extension ratio of 512$\times$, the representation of positions in the initial 4k context window becomes densely packed.

To address this, we conduct an additional evolutionary search on the augmented LLM, specifically targeting the adjustment of RoPE rescale factors for relatively short context lengths (such as 4k and 8k). For these shorter contexts, we decrease the upper bound for the searched $\lambda$ since less positional interpolation is necessary. In the inference phase, the LLM adaptively modifies the relevant RoPE rescale factors.

 \section{Experiments}

\textbf{Evaluation Tasks and Models}. Our experiments integrate LongRoPE with LLaMA2-7B and Mistral-7B, and assess outcomes across three dimensions: (1) perplexity levels for long-context LLMs when processing lengthy texts; (2) the Passkey retrieval challenge, designed to assess how effectively a model can locate a straightforward passkey among a large volume of unrelated content; and (3) performance on conventional LLM benchmarks, using a fixed short context window of 4096 tokens.

\textbf{Fine-tuning}. In the case of LLaMA2, we apply a learning rate of 2e-5 following a linear decay schedule and adopt a global batch size of 32. Fine-tuning is conducted for 400 iterations on the Redpajama~\cite{redpajama} dataset, which is segmented into 128k-length chunks, each delimited with BOS and EOS tokens. Subsequently, starting from the resulting checkpoint, we perform an additional 600 training steps to extend the context window to 256k. Training at the 128k context window is carried out on 8 A100 GPUs using the distributed training platform~\cite{cube}, whereas increasing the context to 256k necessitates 16 A100 GPUs. For Mistral, we utilize a fixed learning rate of 1e-6 with a global batch size set at 64. Both 128k and 256k context versions are fine-tuned in accordance with the approach described in YaRN~\cite{yarn}, using 400 steps and sequence lengths of 16k on the Together Computer's Long-Data Collections~\cite{mistraldata}. This training is executed with 4 A100 GPUs.

\textbf{Search}. For target window sizes up to 256k, the following parameters are employed: $P$=64, $N_1$=$N_2$=16, $p$=0.3, and $\mathcal{T}$=40. In each generation, the top 32 candidates are chosen for mutation and crossover operations. Perplexity evaluation utilizes 5 randomly selected samples from the PG19 validation set, ensuring that each sample meets the target context length. For window sizes exceeding 512k, the values for population size, mutation count, and crossover count are reduced by half. Perplexity in this regime is computed using 3 randomly chosen validation samples from the Pile-Books3~\cite{pile} dataset.

\textbf{Baselines}. To achieve a context window of 2048k, we fine-tuned models on context lengths of 128k and 256k. This results in LongRoPE-2048k (ft=128k) for LLaMA2 and LongRoPE-2048k (ft=256k) for Mistral. For comparison, we evaluate these four models against leading baselines for large context window adaptation, focusing on publicly available LLMs that have undergone positional interpolation-based fine-tuning with PI, NTK, and YaRN methods. The baseline models under consideration include Together-32k~\cite{together}, Code LLaMA~\cite{codellama}, LongLoRA-full-FT-100k~\cite{longlora}, as well as YaRN-LLaMA and YaRN-Mistral~\cite{yarn}.

\subsection{Main Results}

\label{sec:mainresult}
\textbf{Language modeling for sequences up to 256k tokens}. Our initial analysis benchmarks our approach against leading extended LLMs on tasks requiring a context window of 256k tokens. To underscore the method’s versatility, we conduct experiments on the test splits of both Proof-pile~\cite{pg19} and PG19~\cite{pile}. Perplexity is measured at different context lengths by applying a sliding window of size 256. For the PG19 dataset, evaluation utilizes its complete test set comprising 100 documents. In the case of Proof-pile, consistent with YaRN~\cite{yarn}, we draw 10 samples at random, each no shorter than 128k tokens.

Table~\ref{tbl:proof-pile} and Table~\ref{tbl:pg19} present a comparison of the perplexity scores for LLaMA2 and Mistral models enhanced with various interpolation techniques, evaluated on the Proof-pile and PG19 datasets, respectively. Two principal findings emerge: \textbf{(1)} our extended models consistently achieve lower perplexity as the evaluation length increases from 4k up to 256k, indicating their improved capacity to utilize extended context. \textbf{(2)} Despite operating within a context window 16$\times$ longer than standard, which usually complicates the preservation of performance at shorter sequence lengths, our LongRoPE-2048k models surpass state-of-the-art baselines within a 256k context window.

\textbf{Long sequence language modeling beyond 2000k}. To assess performance on very large documents, we utilize the Books3~\cite{pile} dataset. For efficient evaluation, a random sample of 20 books—each longer than 2048k—is chosen, applying a sliding window of size 256k.

As indicated in Table~\ref{tbl:books3}, LongRoPE is able to expand the context window of both LLaMA2-7B and Mistral-7B models up to 2048k tokens, while still attaining perplexity scores that are on par with, or superior to, those of existing baselines for shorter sequence ranges (8k–128k). Substantial differences in performance are apparent between LLaMA2 and Mistral when evaluated at 2048k. Specifically, Mistral exhibits stronger results than baselines on shorter contexts, but its perplexity rises above 7 when evaluated beyond 256k. In the case of LLaMA2, perplexity decreases as the context length increases, in line with expectations, although there are slight upticks at 1024k and 2048k. Additionally, for LLaMA2, LongRoPE-2048k delivers improved results when fine-tuned with a 256k length as opposed to 128k, attributable to a smaller secondary extension ratio (namely, 8$\times$ compared to 16$\times$). Conversely, Mistral shows better performance when fine-tuned using a 128k window. This difference stems mainly from adopting YaRN’s configuration for Mistral, wherein a 16k training length is used during fine-tuning for both 128k and 256k windows; this constrains Mistral’s capacity for additional context window extension post fine-tuning.

\textbf{Passkey retrieval}. Here, we examine how the context window size influences performance on generation tasks. For this, we adopt the synthetic passkey retrieval benchmark introduced by~\cite{passkey}. In this setting, a model must extract a randomly generated five-digit passkey embedded somewhere within a lengthy document. The full prompt template is described in the appendix. We conduct ten rounds of this retrieval task, each time inserting the passkey at a randomly selected position, sampled uniformly throughout the evaluation context window.

Fig.~\ref{fig:passkey} presents a comparison of retrieval accuracy against various baselines. The accuracy of current models falls off sharply to zero when the sequence length exceeds 128k tokens. In stark contrast, LongRoPE-LLaMA2-2048k (ft=256k) demonstrates robust performance in the highly demanding scenario of retrieving a passkey from up to a million tokens, consistently achieving retrieval accuracy of at least 90\% across the entire range from 4k to 2048k tokens. Similarly, LongRoPE-Mistral-2048k (ft=128k) sustains perfect accuracy up to 1800k tokens, only decreasing to 60\% at 2048k. This behavior is consistent with the results in Table~\ref{tbl:books3}, which indicate a moderate increase in perplexity at the 2048k token length.

\textbf{Evaluation on standard benchmarks in the original context window}. The LongRoPE-2048k models are assessed on the initial context window through the Hugging Face Open LLM Leaderboard~\cite{huggingfaceboard} under both zero-shot and few-shot evaluation protocols. Specifically, we utilize 25-shot ARC-Challenge~\cite{allenai:arc}, 10-shot HellaSwag~\cite{zellers2019hellaswag}, 5-shot MMLU~\cite{mmlu}, as well as 0-shot TruthfulQA~\cite{lin2021truthfulqa}.

As indicated in Table~\ref{tbl:commonsense}, our models perform on par with results achieved on the original benchmark, which was intended for shorter context windows, and surpass the original Mistral by +0.5\% on TruthfulQA. Although LongRoPE-LLaMA2-2048k, which was fine-tuned using a 256k context window, experiences a modest decrease in performance, it continues to deliver results that are generally acceptable across the evaluated tasks.

\subsection{Ablation Results}

\textbf{Effectiveness of the second positional interpolation}. Within our progressive extension approach, we employ the search algorithm to perform an additional round of non-uniform positional interpolation on the previously fine-tuned extended LLMs. To assess this method, we conduct experiments on the LLaMA2-256k model that we have fine-tuned. This model is then further scaled up to 512k, 1024k, and 2048k contexts via PI and YaRN. As indicated in Table~\ref{tbl:secondsearch}, our method of non-uniform positional interpolation maintains stable perplexity across these larger context sizes. By comparison, perplexity levels rise rapidly for PI and YaRN as the context length extension increases.

\textbf{Recovery performance at reduced context lengths.} To address reductions in model accuracy when using shorter context lengths, we fine-tune the RoPE parameters for LongRoPE-2048k with our search approach. In this process, we intentionally set lower upper bounds on the scale factors during the search, aiming to limit interpolation for shorter contexts such as 4k and 8k. Table~\ref{tbl:shortperformance} presents perplexity measurements for LongRoPE-LLaMA2-2048k evaluated on the Proof-pile at 4k and 8k, alongside mean benchmark scores for LLM tasks. The data indicate substantial gains in performance at these shorter context lengths.

\textbf{Analysis on the two forms of non-uniformities}. To disentangle the effects of the two types of non-uniformities, we conducted ablation studies to assess their individual contributions to model performance. We designed two experimental settings: (i) we scaled LLaMA2-7B to sequence lengths of 16k and 32k using various strategies—PI, optimizing solely the RoPE dimension, and optimizing both non-uniformity aspects; (ii) we took our LLaMA2 model fine-tuned on sequences of 256k tokens and further extended it to 2048k using analogous methods. All evaluations of perplexity were carried out without any additional fine-tuning. As presented in Table~\ref{tbl:searchdimension}, introducing non-uniformity in the RoPE dimension yields a pronounced improvement over PI’s linear interpolation in terms of reduced perplexity. Non-uniformity in token position produces noticeable gains for 16k and 32k lengths, though these benefits are absent at 2048k—likely because the input sequence is exceptionally long. In scenarios where only the earliest tokens are preserved without interpolation, the technique loses effectiveness; we defer a deeper investigation into this phenomenon for future work.

\section{Related Works}

Beyond position interpolation methods, this section reviews related works employing alternative strategies.

\textbf{Retrieval-based} approaches rely on external memory components to store distant contextual information and utilize retrieval mechanisms to access pertinent documents during inference~\cite{longllama,wang2023augmenting,borgeaud2022improving}. These techniques often require substantial and explicit architectural changes to large language models. By comparison, our method is notably more efficient, involving only minimal changes to positional embeddings. Furthermore, our approach is capable of addressing a broader array of long-context tasks, including long document summarization and few-shot learning, without being restricted solely to retrieval-based scenarios.

\textbf{Attention-based context window extensions}. In addition to positional embedding interpolation, certain studies have managed to extend the input context using the standard LLM context window size by altering the attention mechanisms~\cite{han2023lminfinite, streamingllm,parallelwindow}. The principal concept involves addressing the issue of attention explosion introduced by additional positions through the design of innovative attention masks. Such approaches can be viewed as complementary to methods based on positional interpolation.

\textbf{Fine-tuning based approaches} address the challenge of adapting pre-trained LLMs with altered position embeddings to handle longer contexts. Studies such as Code LLaMA~\cite{codellama}, LLaMA2 Long~\cite{xiong2023effective}, and ScaledRoPE~\cite{liu2023scaling} adopt the strategy of selecting an extremely large base for RoPE, followed by fine-tuning on the desired sequence length. In contrast, our approach supports a wide range of target lengths and is capable of reaching beyond 2M in sequence length. Recently, due to the significant GPU resource requirements associated with fine-tuning for extended context windows (e.g., above 128k), methods like LongLoRA~\cite{longlora} and PoSE~\cite{zhu2023pose} have been introduced to reduce computational overhead. Our technique is complementary to these resource-efficient fine-tuning strategies.

\section{Conclusion}

In this study, we introduce LongRoPE, a technique that significantly increases the context length of LLMs to an unparalleled 2048k, all while preserving their performance within the original, more limited context range. Our approach leverages two distinct types of non-uniformities in RoPE positional embeddings, discovered through a highly efficient evolutionary search process. This methodology yields dual advantages: it delivers effective initial weights for subsequent fine-tuning and supports an eightfold expansion of the context window even without additional fine-tuning. Furthermore, we develop a progressive extension approach whereby LLMs fine-tuned with a 256k context can be scaled to a 2048k window without necessitating further fine-tuning. Comprehensive experimental results substantiate the efficacy of LongRoPE. We anticipate that LongRoPE-2048k models will facilitate a host of novel applications involving extended context and encourage ongoing exploration in this domain.

\section*{Broader Impacts} The research outlined in this paper aims to contribute to progress in the domain of Machine Learning. While our work may have various implications for society, we do not identify any particular effects that warrant special emphasis at this time.

	\balance

\end{document}